\section{Gönüllü Testi Protokolü ve Bilgilendirilmiş Onam Formu}
\label{EK:Gonullu}

\subsection*{Amaç}

Uygulamanın gerçek kullanım koşullarında (telefon kamerası, farklı kişiler, farklı
ortamlar) işaretleri ne kadar doğru tanıdığını ölçmek. Hedef, tez önerisindeki
\sol gerçek ortam doğruluğu yüzde 75\sag{} ölçütüdür.

\subsection*{Katılımcılar}

Ekip dışından, 18 yaşını doldurmuş gönüllüler. İşaret dili bilmek şart değildir;
her işaretin örnek hareketi uygulamada gösterilir. Katılımcı adı hiçbir yere yazılmaz,
yalnızca bir kod (K01, K02, \ldots) kullanılır.

\subsection*{Düzenek}

\begin{itemize}
  \item Telefon ya da tablet sabit bir yere, göğüs hizasında ve yaklaşık 1--1,5 metre
        uzağa konur. Baş, omuzlar ve iki el görüntüde olmalıdır.
  \item Ön kamera kullanılır. Ortam ışığı not edilir (gün ışığı / yapay ışık).
  \item Uygulama uçak modunda çalıştırılır.
\end{itemize}

\subsection*{Uygulama}

\begin{enumerate}
  \item İşaretler sekmesinde sağ üstteki \sol Gönüllü testi\sag{} açılır, katılımcı
        kodu girilir ve 30 işaret seçilir.
  \item Her kelime için katılımcı örnek hareketi en az iki kez izler, \sol Kaydet\sag{}
        düğmesine basar, işareti yapar ve \sol Bitir\sag{} düğmesine basar.
  \item Kayıt hatalıysa (el görünmedi, çok kısa) bir kez yeniden kaydedilir; tahmin
        çıktıysa \sol Sonraki\sag{} ile devam edilir. Yanlış tahminde yeniden deneme
        yapılmaz.
  \item Test sonunda özet ekranı kaydedilir (ekran görüntüsü).
  \item Kayıtlar bilgisayara alınır ve \texttt{python ml/live\_analysis.py} ile incelenir:\\
        {\small\texttt{adb pull /sdcard/Android/data/tr.edu.erciyes.tid\_app/\\files/gonullu\_testi}}
\end{enumerate}

\subsection*{Saklanan veri}

Her deneme için yalnızca el ve gövde noktalarının koordinatları, hedef kelime,
modelin ilk üç tahmini, güven değerleri ve katılımcı kodu saklanır. Kamera görüntüsü
ve ses kaydedilmez. Veriler yalnızca bu bitirme projesi kapsamında kullanılır ve proje
bitiminde silinir.

\subsection*{Bilgilendirilmiş onam formu}

\noindent Erciyes Üniversitesi Bilgisayar Mühendisliği Bölümü bitirme projesi
\sol MediaPipe ve Makine Öğrenmesi Kullanarak Türk İşaret Dili Tabanlı Gerçek Zamanlı El
Hareketi Tanıma ve Çeviri Sistemi\sag{} kapsamında yapılan teste gönüllü olarak
katılıyorum. Testte telefon kamerası önünde işaret dili hareketleri yapacağımı;
görüntümün kaydedilmeyeceğini, yalnızca el ve gövde noktalarının koordinatlarının
adım yazılmadan saklanacağını; verilerin yalnızca bu proje için kullanılıp proje
sonunda silineceğini ve testi istediğim an bırakabileceğimi biliyorum.

\vspace{1cm}
\noindent
\begin{tabular}{@{}p{0.45\textwidth}p{0.45\textwidth}@{}}
Katılımcı kodu: \dotfill & Tarih: \dotfill \\[8mm]
İmza: \dotfill & Testi yürüten: \dotfill \\
\end{tabular}
